\itemhead{复习题三 · 第10题（教材p.90）}{XJ-B1-C03-RV-P090-Y1501}
\label{XJ-B1-C03-RV-P090-Y1501}
{\footnotesize 湘教版教材p.90 / PDF第97页 / 复习题三 / 第10题}
\field{原题}{已知函数 $f(x)=(x+1)(x+a)$ 为偶函数，求 $a$ 的值。}
\field{参考答案}{$a=-1$。展开后一次项系数 $a+1$ 必为0；此时函数为 $x^2-1$。}
{\small \field{关联知识}{F08 奇偶性、对称与函数性质应用（教材帮p.113、114、115 / PDF122、123、124）}
\field{方法与易错点}{检查定义域，再代入 $-x$；求另一侧解析式时同时写出 $-x$ 所在区间。 易错：只看公式不看定义域；用一对点证明奇偶；负半轴代换遗漏最外层负号。}
\field{答案依据}{答案册相关题组：PDF第24、25、26、27页（册内页23、24、25、26）；按当前题库原题逐项重算，不按旧版题号机械对套。}}
